\section{Discussion}
\label{sec:discussion}

\subsection{Lessons}
\label{sec:lessons}


\textbf{Pre-environment scoring surfaces a failure class that post-environment benchmarks structurally cannot see.} 69.8\% of our runs produce simulated reproductions which never execute the real firmware binary. This failure class is invisible under post-environment scoring. The lesson for benchmark design is concrete: record full artifacts and trajectories, require artifact provenance for every claimed success. Without these, an agent can score as successful by emulating the vulnerability while never touching the target.

\noindent{\textbf{On from-scratch reproduction, behavior is the decisive factor.}} \texttt{glm-5.2} outperforms \texttt{claude-sonnet-4-6} and \texttt{gpt-5.5} despite being a smaller model, because it persists longer at firmware acquisition, plans better fallbacks, and resists simulation substitution. The observation that stronger models do not necessarily substitute less indicates that simulation is driven by task difficulty and behavioral defaults, not by model weakness.

%%\textbf{The rehosting ceiling is a specific, %%identifiable blocker---not a general capability %%wall.} The architecture-dependent ceiling we %%observe---same-architecture (ARM64) devices are %%reproducible via \texttt{chroot} while %%cross-architecture (MIPS, ARM32) devices almost %%universally fail at P5---localizes the blocker %%precisely. The failing sub-tasks are not vague: %%NVRAM initialization, dynamic-loader and %%library-path resolution under uClibc/musl, and %%network-bridge setup are exactly the rehosting %%challenges documented in the firmware-analysis %%literature~\cite{challenges,pandawan}. The %%implication is that progress on \sys{} will come %%not from larger general- purpose models but from %%agents that can reliably perform the %%configuration steps a human analyst performs by %%hand, or that can invoke and drive specialized %%rehosting toolchains~\cite{firmadyne,firmae,%%greenhouse}. A post-environment benchmark would %%never surface this ceiling, because the service %%is already running; \sys{} exposes it precisely %%because it scores the rehosting phase directly.

\subsection{Limitations}
\label{sec:limitations}

\textbf{Dataset scope.} \sys{} contains 30 manually curated CVEs balanced equally across buffer overflows, command injections, and authentication bypasses. This balance is a deliberate design choice that enables controlled comparison across classes requiring distinct validation signals, but it does not mirror the natural distribution of IoT weaknesses.

%%, which is dominated by command injection and %%memory-safety bugs~\cite{cwe2025kev}. The corpus %%is also confined to Linux-based firmware for %%SOHO-class devices; real-time-OS, %%industrial-control, and smart-home ecosystems %%with proprietary stacks are not represented. %%Thirty CVEs are sufficient to expose a dominant, %%cross-model failure mode---simulation %%substitution appears in every model and every %%class---but they bound the statistical power of %%per-class and per-vendor comparisons, so we %%report aggregate trends and qualitative class %%effects rather than class-level significance.

\noindent{\textbf{Scoring Skill.}} \sys{} use the scoring skill which is a rubric-constrained evaluator rather than a deterministic exploit oracle: it audits artifact, but a stratified human audit, not automated execution, validates the non-zero P5/P6 claims on which the headline success rates rest.

\noindent{\textbf{Single Agent Harness.}} To isolate model behavior from harness variance, we run all five models through a single unified harness with an identical prompt, tool set, and control loop. Agent frameworks differ in their tool interfaces, context management, retry and recovery policies, and planning loops; alternative harnesses such as Claude Code~\cite{claude_code2026} or custom multi-agent pipelines may shift both the overall scores and the simulation-substitution rate, particularly if a harness injects explicit real-target verification steps between phases. Likewise, our five models span a range of providers and sizes but do not exhaust the space. Extending \sys{} to multiple harnesses and a wider model and agent pool would test whether the behavioral factors we identify are stable across harnesses or are artifacts of a single configuration. We accordingly view the single-harness, five-model study as a first controlled slice of this evaluation space, not as a ceiling on agent capability.